\chapter{Nonrelativistic results}\label{ch:nonrel}

The first stage of the work used the Hermite--Fourier Vlasov solver of the parent code with the Maxwell--Proca coupling added. The model is nonrelativistic, deterministic and noise-free, and its energy ledger closes to solver tolerance, which makes it an unusually clean instrument for asking how much of the HHS saturation can be a property of the plasma physics as such. This chapter collects what the instrument established: where it can be trusted, what the transfer looks like inside those limits, how a finite dark-photon reservoir, a finite coherence time, collisions and a density gradient modify it, and why all of it, however carefully certified, answers a question about a model that is not the one HHS simulate. The last point is the reason the chapter ends where it does, with the quiver speeds of Section~\ref{sec:c9-scope}, and the reason that Chapter~\ref{ch:rel} exists.

Throughout, $\vq$ is the instantaneous drive strength defined in Chapter~\ref{ch:hhs}, the box is $L=40\,c/\wpe$, $m_i/m_e=1836$, $T_i=T_e$, $v_{te}^2=10^{-3}c^2$ and the drive frequency is the Langmuir frequency $\omega_L=\wpe\sqrt{1+m_e/m_i}$ unless stated. The initial state carries a deterministic density seed of relative amplitude $5\times10^{-4}$ on the lowest Fourier modes, and not the shot noise of a particle code.

\section{Certified windows at strong drive}\label{sec:c9-windows}

\subsection{The truncation instability and what a window is}

The Hermite--Fourier discretization expands the electron distribution in asymmetrically weighted Hermite functions in velocity and in Fourier modes in $x$. In a spatially uniform field the operator $-v\partial_x+(q/m)E\partial_v$ reduces to a banded matrix whose action on the first two Hermite coefficients is exact, which is why the homogeneous response is reproduced to $10^{-9}$ for $\wpe t\le1000$ even at quiver speeds of $15\,v_{te}$. In a non-uniform field the truncated operator is not anti-self-adjoint. Its spectrum for a frozen field has growing modes in the top third of the Hermite ladder, at the largest retained wavenumber, with a rate that increases with the Fourier cut-off and with the number of Hermite functions $N_n$: at $K=5$, $N_n=128$ and a field amplitude $10^{-2}$ the largest growth rate is $1.6$, while the same Galerkin operator written in symmetric Hermite functions is neutral to round-off ($8\times10^{-15}$; record \texttt{studies/highk/eig.json}). This is a defect of the basis and not physics, and it must not be confused with the physical oscillating-two-stream instability, which the Hermite model reproduces within $1.3\,\%$ of Silin's kinetic dipole-pump relation~\cite{Silin1965} over 22 scan points and which drives the genuine cascade to higher $k$.

A field-scaled closure damps the top of the ladder at a rate proportional to the instantaneous field, and it is exact in the Lawson factor so that it adds no stiffness to the time integration. It is the only remedy we found. The remaining question, which governs everything below, is how long a run can be trusted. We call a \emph{certified window} the time up to which four declared checks hold. Refinement pairs in $N_x$, $N_n$ and closure strength must agree to $10\,\%$ in the electron kinetic-energy change $\Delta K_e$ and in $\Wext$. The electron distribution must stay positive, which we monitor through the positivity of the Hermite-zero kinetic energy and of the $k=0$ temperature. The solver must reach the time. The energy ledger must close. The window is the smallest of the four times. It is a lower bound on the time to which the Hermite results are reliable, and not a statement about the plasma. A rejection of a window by one of the tests never means that the physics changes there, only that the numerical instrument has stopped being demonstrably right.

\subsection{Windows and references}

The windows were established in three campaigns. The first, the $N_x$-convergence study (\texttt{nxconv}), compared Hermite runs at $N_x=32,64,128$ with a masked spectral grid code that integrates the Vlasov equation in $x$--$v$ without any Hermite truncation, at $\vq=0.03$ and $0.1$. At $\vq=0.03$ the Hermite code at closure strength $c=4$ agrees with the grid in $\Delta K_e$ and $\Wext$ to $\wpe t\simeq960$--$1000$ for $N_x\ge32$, with field spectra agreeing to $720$--$770$. At $\vq=0.1$ the window closes near $\wpe t\simeq600$: the pair $N_n=64$ against $128$ separates at $598$, the pairs at different $N_x$ at $620$ for $c=1$ and to $760$ for $c=4$, and the field spectrum agrees only to $555$. Increasing $N_x$ or $N_n$ ends the run sooner, which is the signature of an instability and not of under-resolution. The second campaign (\texttt{gridref}) asked whether the grid reference could itself reach $1000$ at $\vq=0.1$. With the velocity box enlarged to $48$ and then $64\,v_{te}$ the first time at which the edge fraction exceeds $10^{-6}$ moves only from $701.5$ to $754.5$, because the tail keeps accelerating beyond $64\,v_{te}$; the first appearance of $f<-10^{-3}$ is at $645.5$ for every $N_x=64$ run whatever the box and velocity resolution, and at $498$ for $N_x=32$, so spatial resolution sets it. An earlier statement in the repository that the box edge limits the grid near $650$ was wrong and was corrected. The grid reference at $\vq=0.1$ is certified for $\Delta K_e$ and $\Wext$ to about $750$ and for the spectrum to about $640$. The third campaign (\texttt{highk}) established the instability itself and the closure, and ruled out aliasing (a padded convolution equals the $2/3$-masked right-hand side to $10^{-16}$), the step size (a fixed step fails at the same time) and the segment remaps (without remaps the run fails at $472$).

Three attempts to extend the windows failed under declared checks, and they belong in this chapter because they bound what can be said. A GPU campaign raised the resolution to $N_x\in\{64,128\}$ and $N_n\in\{128,256\}$ at $\vq=0.1$ and $0.05$; every run ended at the step budget inside the instability, between $\wpe t=328$ and $595$ at $\vq=0.1$ and between $440$ and $751.5$ at $\vq=0.05$, and the pairs agreed to the end of the shorter run, so the failure is of the route and not of the agreement. A closure-strength study derived the frozen-field growth rate as a function of closure strength and found that no strength up to $1024$ brings it below $10^{-3}$ at all 39 sampled points with a peak field of $4\times10^{-3}$; the floor grows as $K^{0.15}N_n^{0.30}E^{1.26}$, so a resolution-consistent closure strength does not exist for this closure family. A stronger fixed closure did extend the $\vq=0.1$ window from $595$ to $713$--$757$, but the runs end by positivity loss at a time that coincides with the limit of the grid reference, so the cause cannot be attributed to the Hermite floor alone. Finally three alternative bases (the Galerkin form of Dai~\cite{Dai2026}, the symmetric-weight basis, and the anti-symmetric square-root form of Issan et al.~\cite{Issan2024a,Issan2025}) are linearly neutral and accurate against the kinetic dispersion relation, but fail their declared nonlinear gates at $\vq=0.1$: the symmetric basis blows up at $637$--$738$, and the square-root form departs from the grid at $391$ ($N_x=16$) or $498$ ($N_x=32$), earlier than the symmetric one and independently of $N_n$ and the time step. The asymmetric basis with closure remains. The status is summarized by the windows of Table~\ref{tab:c9-win}, which are the numbers used in the rest of this chapter.

\begin{table}[t]
\centering
\caption{Certified windows for the strong-drive set-up, closure strength 64 integrated exactly, members $(N_x,N_n)=(32,64),(64,64),(32,128),(64,128)$ and $(32,64)$ at $c=32$ (records \texttt{studies/r1nl/stage1\_summary.json}). The last column gives the time to which the spectral-grid reference at $N_x=32$ and 64 agrees with the Hermite coarse member (reported, not part of the window). At $\vq=0.1$ no window beyond about 600 exists (Section~\ref{sec:c9-windows}).}\label{tab:c9-win}
\small
\begin{tabular}{@{}cccc@{}}\toprule
$\vq$ & window $\wpe t$ & limit & grid agrees to\\\midrule
0.03 & 1000 & cap of the study & 1000\\
0.05 & 892 & positivity, $N_x=64$ members & 979 / 962\\
0.07 & 791.5 & positivity, $N_x=64$ members & 870 / 863\\
0.1 & $\sim600$ & truncation instability & $\sim750$ (reference only)\\\bottomrule
\end{tabular}
\end{table}

\section{Linear transfer, finite reservoir and finite coherence}\label{sec:c9-r1}

\subsection{The transfer is linear in the window}

Inside the windows of Table~\ref{tab:c9-win} the work follows the linear resonant law. At the strongest certified drive, $\vq=0.07$, the maximum of $|\Wext/\Wlin-1|$ over four resolutions is $6.8\times10^{-3}$ and the four members agree on it to $2\times10^{-4}$; at the window end $\Wext=381\,nT_eV$ against $\Wlin$ within the same fraction. The mean-field frequency shift measured as the slope of the lock-in phase difference is below $6.5\times10^{-6}\wpe$, against a declared resolution threshold of $10^{-3}\omega_L$, so it is not resolved. No saturation and no detuning are resolved, and the declared departure rule, a $10\,\%$ departure in the coarsest member confirmed at $5\,\%$ in all refinements, is not met. The ledger closes to $10^{-9}$ over the 33 runs of the study (declared bound $10^{-6}$).

The energy goes into the coherent quiver motion. At the window end of $\vq=0.07$, $370$ of the $381\,nT_eV$ of work is coherent $k=0$ electron kinetic energy and the ions receive $0.05\,\%$ of it. With the mean flow removed the thermal gain over the trailing $200\,\wpe^{-1}$ is $0.14$--$0.24\,nT_eV$ (that is $0.28$--$0.48\,\Eth$ by the factor of two of Chapter~\ref{ch:hhs}), and the HHS measure computed for the three drives is $2.5\,\Eth$ at $\vq=0.03$ (at $\wpe t=1000$), $4.5\,\Eth$ at $0.05$ and $5\,\Eth$ at $0.07$ at the window ends, against $0.1$, $0.28$ and $0.54\,\Eth$ at $\wpe t=500$. These are small numbers in comparison with the $50\,\Eth$ of HHS and with the work, and the oscillation behind them is informative: the electron motion at $k\ne0$ sloshes between wave kinetic and field energy with a period of about $67\,\wpe^{-1}$ near the onset, so that a single time sample is not a heating measure, and only the average is. The physical oscillating-two-stream instability turns on at $\wpe t\simeq780$--$860$, depending on $N_x$, and the fluctuating field $\sum_{k\ge1}|E_k|^2$ grows appreciably only then. \emph{This is the limitation of the bound}: it says that no saturation appears before the cascade starts, and nothing about what the cascade does afterwards, since the run ends by positivity loss soon after.

\subsection{The seed back-reaction}

A small, systematic deficit remains in $\Wext/\Wlin-1$, and understanding it took a hypothesis and a rejection. At the three drives the deficit is the same to three digits, it grows as $-10^{-8}t^2$ until the instability stage, and it reaches $-1.07$, $-0.88$ and $-0.68\,\%$ at the window ends $1000$, $892$ and $791.5$. The deficit is therefore independent of $\vq$ in the range tested, which is surprising for a nonlinear effect of the drive. Two hypotheses were declared. The first was that it is a truncation effect of the seedless homogeneous system, and it was rejected: the unseeded deficit, $2.3\times10^{-9}$ at $\wpe t=400$, is only $1.4\times10^{-6}$ of the seeded one. The second was that it is a back-reaction of the seed. A seed-scale test multiplied the initial density fluctuations by $s=1$, $0.1$ and $0$ at $\vq=0.07$ and found $-1.66\times10^{-3}$, $-1.66\times10^{-5}$ and $2.3\times10^{-9}$ at $\wpe t=400$; the ratio of the $s=0.1$ to the $s=1$ deficit is $0.0100$, equal to $s^2$.

The origin of this term is visible in the mean-field identity. The $k=0$ momentum equation is exact, $\dot{\bar J}=\omega_L^2(\bar E+E_0\cos\omega t)+\mathcal N$ with $\mathcal N\propto\langle\delta n\,\delta E\rangle$, and the quiver advects the seed modulation, so that density and field fluctuations of the seed, each linear in the seed amplitude, multiply to a $k=0$ contribution proportional to the seed amplitude squared. It is a response to the imposed fluctuation level, not to the pump, and a physical fluctuation level different from $5\times10^{-4}$ would rescale it. The distinction matters for what follows: the same $s^2$ deficit reappears in the collisional study as the failure of the detuned low-collisionality cases (Section~\ref{sec:c9-r2}), and in the comparison with the PIC it is the first hint that the fluctuation level, and not the drive, controls how far a calculation departs from the linear law.

\subsection{Finite reservoir}

In the self-consistent mode the dark photon is a free $k=0$ Proca mode with $\OmD=\omega_L$ whose energy $U_{D0}$ is finite and is depleted by the transfer. The reservoir is parametrized by $u=U_{D0}/(nT_eV)$ and the coupling $\eta$ is chosen such that the initial effective force matches the prescribed drive, $\eta=(\vq)/\sqrt{2u}$. Because the mobile-ion plasma and the dark mode form a pair of coupled oscillators at the same frequency, the exact linear dynamics is a four-dimensional system for $(\bar E,\bar J,E_D,A_D)$ whose matrix exponential agrees with an ODE solution to $7\times10^{-11}$. The kinetic dark energy $U_D(t)$ follows this linear prediction to within $0.039\,U_{D0}$ in six cases ($u=0.1,1,10$ at $\vq=0.01$ and $0.03$, windows $1000$ and $940$), to $0.23\,\%$ and $0.81\,\%$ of $U_{D0}$ at $\vq=0.07$ for $u=1$ and $10$, and the beat times equal $\pi/\eta$ (for example $140$, $444$, $46$, $148$ and $468$ in the five cases that reach them; the $u=10$, $\vq=0.01$ case does not reach its first minimum in $1000$). The accumulated work $W_D$ follows the matched prescribed drive until the split time $t_{\rm split}=50.5$ and $72$ for $u=1$ and $10$ at $\vq=0.07$, equal to the linear values, and remains within $1.0$ and $2.7\,\%$ of the linear prediction while it exceeds $20\,\%$ of its running maximum. A finite reservoir is thus described by the linear theory over the windows; no nonlinear reservoir response is resolved.

\subsection{Finite coherence}

Real dark-photon dark matter is not perfectly coherent. We modelled a finite coherence time by a prescribed drive $E_0\cos(\omega_Lt+\Phi(t))$ with $\Phi$ a phase random walk, giving a Lorentzian line of half-width $1/\tau_c$ and no phase resets, with six kinetic realizations per $\tau_c$ against 400 realizations of the linear theory. The first version realized a phase variance of $dt/\tau_c$ instead of the declared $2\,dt/\tau_c$, so its coherence times were twice nominal; the records were kept and labelled as such and a corrected version was declared and run before any new result was inspected (the product rate times $\tau$ in the corrected version is $1.99$). The ensemble ratio of the kinetic to the linear same-drive work is $1.000\pm0.007$, $1.089\pm0.064$ and $0.998\pm0.003$ for $\tau_c\wpe=30$, $100$ and $300$ at $\vq=0.01$, window $1000$, where the $\tau_c=100$ value is $1.4$ standard errors above unity and consistent with sampling; at $\vq=0.07$ the ratios are $1.0050\pm0.0058$ and $1.0000\pm0.0038$ for $\tau_c\wpe=100$, $300$, which bounds the ensemble nonlinear correction by $1.7$ and $0.8\,\%$. Individual realizations deviate more where the linear work is itself small. The late work rate saturates at about $E_0^2\tau_c\omega_L^2/4$ rather than growing, but this is the linear response to a broadened line and not a nonlinear effect: for $\tau_c\wpe=30$, $100$ and $300$ the derived slopes are $8.1\times10^{-7}$, $2.6\times10^{-6}$ and $7.2\times10^{-6}$ against $2.25\times10^{-5}$ at the end of the coherent case, where the work is still growing linearly in time.

The bounds hold for one box, one ion mass, $T_i=T_e$, drives to $\vq=0.07$, one seed amplitude, a Lorentzian line only, and reservoir and coherence not combined in one run. They do not address the inhomogeneous dark fields or the Gaussian line.

\section{Collisional absorption and the Drude law}\label{sec:c9-r2}

\subsection{The linear drag law}

The cosmological case differs from the collisionless one because the energy is dissipated: collisions with ions bring a steady absorption that does not saturate and is the quantity on which the bounds are in fact built. With a collision operator that conserves momentum and energy between species, the relative velocity of electrons and ions decays at the Morse rate, and for a spatially uniform driven plasma the cold equations close,
\begin{equation}
\dot{\bar E}=-\bar J,\qquad
\dot{\bar J}=\omega_L^2\bigl(\bar E+E_0\cos\omega t\bigr)-\nu_{\rm eff}\bar J,\qquad
\nu_{\rm eff}=\frac{\nu_{ei}}{2}\Bigl(1+\frac{m_e}{m_i}\Bigr).
\label{eq:c9-drude}
\end{equation}
The factor $\tfrac12(1+m_e/m_i)$ is the one the Dougherty operator gives for the momentum friction in the two conventions we use, the moment-conserving multispecies form of Francisquez et al.~\cite{Francisquez2022} (LBO-EM, parent PR 71) and the same operator with the same operator with $\beta=0$ (LBO-G, PR 73), which differ only in the temperature-difference exchange. Eliminating $\bar J$ gives a damped oscillator, $\ddot{\bar E}+\nu_{\rm eff}\dot{\bar E}+\omega_L^2\bar E=-\omega_L^2E_0\cos\omega t$. The periodic solution is $\bar E={\rm Re}[\hat Ee^{i\omega t}]$ with $\hat E=-\omega_L^2E_0/(\omega_L^2-\omega^2+i\nu_{\rm eff}\omega)$. The mean power absorbed from the drive is $P=\langle\bar JE_0\cos\omega t\rangle=-\langle\dot{\bar E}E_0\cos\omega t\rangle=\tfrac12E_0\omega\,{\rm Im}\,\hat E$, which evaluates to
\begin{equation}
P_{\rm ss}=\frac{\omega_L^2\,\nu_{\rm eff}\,\omega^2E_0^2}{2\bigl[(\omega_L^2-\omega^2)^2+\nu_{\rm eff}^2\omega^2\bigr]},\qquad
P_{\rm ss}^{\rm res}=\frac{\omega_L^2E_0^2}{2\nu_{\rm eff}}.
\label{eq:c9-pss}
\end{equation}
Two features deserve notice. The resonant absorbed power is inversely proportional to the collision rate, so that a weaker friction absorbs \emph{more} in steady state, but only after the amplitude transient, which decays as $\exp(-\nu_{\rm eff}t/2)$, has died; the one-over-$e$ time is $4/\nu_{ei}$ and for $\nu_{ei}=10^{-3}\wpe$ it is four times the length of any run we can afford. Away from resonance the power is proportional to $\nu_{\rm eff}$, as in the usual Drude picture of off-resonant absorption. Equation~\eqref{eq:c9-pss} is the cold Drude response of Dubovsky and Hern\'andez-Chifflet~\cite{DubovskyHC2015}, who write the friction as $-m_e\nu(u_e-u_i)$ with a Spitzer-like $\nu$, and agrees with the non-resonant absorption rate of McDermott and Witte~\cite{McDermottWitte2020}; their $\nu$ corresponds to $\nu_{\rm eff}$ here. In our model $\nu_{ei}$ is a constant. The Spitzer scaling $\nu\propto T_e^{-3/2}$ is absent, and we return to this in Chapter~\ref{ch:discussion}.

\subsection{Kinetic against linear}

The test is whether the full kinetic solution, with its seed and its heating, follows Eq.~\eqref{eq:c9-drude}. The grid had $\vq=0.01$, $0.03$, $0.05$, $\nu_{ei}\wpe^{-1}=10^{-3}$, $3\times10^{-3}$, $10^{-2}$, $3\times10^{-2}$ and $\omega/\wpe=\omega_L$, $0.9$, $1.1$: 192 production runs, including $N_x=64$ and $N_n=128$ refinement members and unseeded controls, with every check declared in \texttt{studies/r2\_gates.py} before the first run. All LBO-EM cases reach $\wpe t=1000$, refinement pairs agree to the end, and the late-rate spread among refinement members is at most $2.2\times10^{-4}$. The result is that the kinetic absorption equals the linear drag law wherever the law is testable. At $\nu_{ei}=3\times10^{-2}$, where the steady state is reached, the late rate is within $0.4\,\%$ of $P_{\rm ss}$ for all three frequencies and drives ($P_{\rm kin}/P_{\rm ss}=0.996$--$1.003$), and the steady-state law holds in all 12 cases where the linear transient has decayed to within $10\,\%$ of $P_{\rm ss}$. For the other cases the steady state is not reached in $\wpe t=1000$: the resonant rate at $\nu_{ei}=10^{-3}$, $3\times10^{-3}$ and $10^{-2}$ is only $0.20$, $0.48$ and $0.89$ of $P_{\rm ss}$, and these cases are tested against the exact transient-inclusive solution of Eq.~\eqref{eq:c9-drude}, to which the kinetic rate agrees within $0.2$--$4\,\%$. The gate of $5\,\%$ agreement passes in $30$ of $36$ cases. The six failures are the detuned $\nu_{ei}=10^{-3}$ cases, where the late rate is $9$--$14\,\%$ high independently of $\vq$; the unseeded controls agree with the linear law to $3\times10^{-8}$, and a seed-scale follow-up declared before its runs gave ratios of $0.0105$ and $0.0099$ at $s=0.1$ against $s^2=0.01$. These are the same seed back-reaction as above, large only where the linear work is itself tiny because the drive is detuned and the collisions are weak. The resonant late rate at $\nu_{ei}=10^{-3}$ is $1.5$--$4\,\%$ low, and at $\vq=0.05$, $\nu_{ei}=10^{-3}$ it is $3.7\,\%$ low with an electron temperature $129$ times the initial one; drive and heating are not separable in these data and we do not attribute it.

\subsection{Scaling with drive and collision rate}

Fitting $\ln P=a+p\ln\vq+q\ln\nu_{ei}$ over the twelve cases per frequency gives for the resonant, $0.9$ and $1.1$ drives $p=1.99\pm0.10$, $2.00\pm0.03$ and $2.00\pm0.02$, against exactly 2 in the linear law. The collision exponent follows the transient-inclusive linear law and not a pure power: $q=-0.51$ in the resonant case against $-0.52$ for the linear law (because the steady state is not reached, the pure resonant exponent $-1$ of Eq.~\eqref{eq:c9-pss} is not seen), and $0.75$ and $0.88$ against $0.79$ and $0.90$ for the detuned frequencies. The geometric-mean ratio of kinetic to linear power is $0.985\pm0.003$ in the resonant case and $1.040\pm0.017$, $1.025\pm0.011$ for the detuned frequencies. The domain of the law is therefore $\vq\le0.05$, $\nu_{ei}\in[10^{-3},3\times10^{-2}]\wpe$, $\omega\in\{0.9,\omega_L,1.1\}\wpe$, $t\le1000$, one box, one mass ratio and one seed level.

\subsection{LBO-EM against LBO-G, and heating}

The two operators agree where they should. In $\Wext$ they differ by $2.5\times10^{-6}$ at most, the late-rate ratio by $7\times10^{-6}$, and the electron kinetic energy and temperature by at most $8\times10^{-3}$ over all 36 pair comparisons. They differ where they should too: the ion temperature differs by up to $31\,\%$ and the electron--ion energy exchange is halved (its relative change is $-0.4999$ to $-0.4987$), because the absorption is set by the momentum friction, identical in the two forms, whereas the exchange is the temperature-difference rate that LBO-G reduces. A separate LBO-G study declared a $2\,\%$ agreement gate for $\Wext$, $\Delta K_e$, $T_e$ and $T_i$ and failed it through $T_i$ alone ($-5.3\,\%$ at $\wpe t=750$ and $-9.8\,\%$ at $950$ for $\vq=0.03$, $\nu_{ei}=10^{-2}$); the gate was not amended since $T_i$ is fed by the exchange itself, and the other three quantities change by at most $1.8\times10^{-3}$. Ion heating is small ($T_i/T_i(0)\le2.8$) and the electron-to-ion exchange is $4\times10^{-5}$ to $1.6\times10^{-2}$ of $\Wext$, proportional to $\nu_{ei}$.

The electrons heat strongly: $T_e/T_e(0)$ reaches $231$, and exceeds $90$ ($v_{te}=0.3c$) in the resonant $\vq=0.05$ cases, which is outside the nonrelativistic model; the nonrelativistic window for these cases is shortened to $912$, $742$, $616$ and $738$ for the four collision rates. Heating does not change the rate in this model, because $\nu$ is constant: the late-rate ratio stays within $1\,\%$ of unity while $T_e/T_e(0)$ rises from $7$ to $70$ at $\vq=0.03$, $\nu_{ei}=10^{-2}$, resonant. In a real plasma it would, since $\nu_{ei}\propto T_e^{-3/2}$ falls as the plasma heats and lowers the absorption. A cosmological deposition law needs therefore at least the temperature-dependent Coulomb frequency, radiative and Compton cooling, expansion and the ionization history, relativistic electrons, the $k\neq0$ coupling of a real dark-photon field, and a steady-state measurement at $\nu_{ei}\le10^{-2}$ that needs $\wpe t\gg4/\nu_{ei}$. None of these is in the present study. The checks that failed along the way were a declaration defect in the steady-state self-check (window 502 longer than the record of 400; the corrected check, declared afterwards, passes at $3\times10^{-5}$) and, in the earlier collision pilot, a zero-collision rerun that differed from the old record by $6\times10^{-6}$ against a gate of $10^{-6}$, traced by a declared follow-up to the restart chunk length (240 against 200) and not to the solver change. Both are recorded, and neither modifies a claim.

\section{Inhomogeneous and magnetized conversion}\label{sec:c9-r3}

HHS contrast the resonant case with a Landau--Zener passage, in which the plasma frequency sweeps through the dark-photon mass. In a plasma with a density gradient and a magnetic field the question has a spatial form, and it was the object of the third research study (R3). It uses packets that cross a tanh density ramp through the resonance $\wpe=\OmD$, for $\eta=0.15$ and $0.3$, both polarizations, with a uniform magnetic field absent, along, oblique to and perpendicular to the gradient ($|\Omega_c|=0.3$). Two codes were compared: an independent NumPy cold-fluid Maxwell--Proca time-domain reference (19 fields, pseudo-spectral with $N=4096$ and an eighth-order Dormand--Prince integrator) and Dark-SPECTRAX in 1D3V, with the kinetic equilibrium solved exactly on the discrete first-order row so that its drift is $3\times10^{-16}$. The reference conserves energy to $3.5\times10^{-10}$, converges under refinement to $4.5\times10^{-5}$, and balances the energy flux to $3.4\times10^{-6}$ for transverse cases only when the Proca flux contains the term $\epsilon_0\OmD^2\phi_D\mathbf A_D$; without it the budget is off by $0.02$ to $0.9$ of the incident energy. For transverse packets, 16 runs, the kinetic solution equals the cold reference to $4\times10^{-4}$ with ledger $\le2\times10^{-10}$ of the incident energy.

\subsection{Landau--Zener and the Dykhne exponent}

The standard Landau--Zener formula for two coupled modes whose frequency difference sweeps linearly, written for a spatial ramp, gives a transmitted fraction $P=\exp(-\pi g^2/(\bar k\,|\mathrm d\Delta/\mathrm dx|))$, with $g$ the coupling, $\bar k$ the wavenumber at the crossing and $\Delta$ the local frequency mismatch. For a transverse packet we found, at the end of the sweep, that the transmitted energy is $P_d$, the reflected photon $P_d(1-P_d)$ and the reflected dark photon $(1-P_d)^2$, the classical two-state result. The simple formula was subjected to a declared accuracy check (gate R4) on four chosen cases, and it failed in two of them, by $0.045$ and $0.035$ at $\eta=0.15$, ramp widths $12$ and $8$. The exponent is too large by a factor $1-4.5\,s^2$ with $s=\eta n_c/\bar k^2$ ($n_c$ as defined in the study script \texttt{studies/r3\_followup\_gates.py}): the ratio of the observed to the simple exponent is $0.99$, $0.95$, $0.89$, $0.80$ and $0.61$ at $\eta=0.05$, $0.1$, $0.15$, $0.2$ and $0.3$, so the simple law is accurate to $0.01$ for $s\lesssim0.1$ only. A refinement of the exponent by the complex-turning-point method (Dykhne), declared before it was computed (follow-up G), reproduces all three channels to $0.009$ for $\bar kL_c\ge3$ at all $\eta$, and to $0.025$ for the narrowest ramp. The relative error of the simple law scales as $\eta^2/\bar k^2$, which is not small only if $\eta$ is comparable to the wavenumber of the packet; for non-relativistic dark matter with $\bar k\ll1$ the condition is that $\eta\ll\bar k$, which is far from satisfied by $\eta=0.3$ and was not tested at the realistic regime. A further prediction declared before its run (gate R5) was wrong and is recorded as such. The relevance to HHS is limited by two facts that bound the result in the other direction: the packets are relativistic ($k_0=1$), unlike dark matter, and the amplitude is linear, so the nonlinear suppression HHS assert for the conversion is not tested. The calculation of Beadle, Caputo and Ellis (arXiv:2405.13882), whose abstract we read, solves a second-order boundary-value problem for non-relativistic dark matter in a strongly varying plasma; our time-domain cold solve is the same linear problem at relativistic $k_0$ and is not a reproduction.

\subsection{Longitudinal trapping and the effect of $\mathbf B_0$}

For longitudinal polarization the picture is different. The wave tunnels as $\exp(-2\int\kappa\,\mathrm dx)$ and the transmitted dark fraction equals this WKB estimate (at most $1.4\times10^{-3}$ for ramp widths $\ge3$), but up to $28\,\%$ of the packet energy stays at the cold resonance layer, where phase mixing generates $k\simeq40$ by the end of the run, which needs $N_x\ge2048$. At $N_x=3072$ the kinetic and cold solutions agree to $6\times10^{-4}$ on the propagating channels and the layer fraction is $0.205$ against $0.199$; the temperature range $v_{te}=0.005$--$0.02$ shows no resolved kinetic correction at this resolution. A flux budget check fails at $N_x=512$ for the narrowest ramp ($1.5\times10^{-2}$) and for the resonance cases, and the refinement check fails for longitudinal and oblique packets at $N_x=512$ and $768$ (eight of 15 pass), for the simple reason that the layer is not resolved. Thermal absorption in the layer, which would need $N_x\gtrsim12000$ and $N_n\gg4$, is not done.

A uniform magnetic field leaves the transmitted fraction unchanged (at most $0.004$ for ramp widths $\ge4$ and $0.02$--$0.06$ at width $2$); a field along the gradient, and the ordinary mode, behave exactly as with $\mathbf B_0=0$. With the field oblique or perpendicular, however, the trapped longitudinal energy of the layer ($0.20$) is converted into reflected ordinary photons ($0.26$--$0.27$), and $0.05$--$0.10$ of the transverse packets goes into a hybrid layer, with the reflected X-mode photon dropping from $0.25$ to $0.05$. This is a statement of where the energy goes in the cold-plasma limit and with immobile ions, in one dimension; it says nothing about thermal absorption at the layer, about mobile ions, or about two dimensions.

\section{What these results do and do not say}\label{sec:c9-scope}

All of the statements above are about a nonrelativistic model. The coherent quiver speed of the electrons implied by the work, $v=\sqrt{2\Wext/(nm_e)}$, is $0.47c$ at the end of the $\vq=0.03$ window, $0.70c$ at the end of the $0.05$ window and $0.87c$ at the end of the $0.07$ window. These can be understood from the resonant growth $v_m\simeq\vq\,v_{te}\,\wpe t/2$, which with $v_{te}=0.0316c$ gives $9.5\times10^{-4}c\cdot t/2=0.47c$ at $t=1000$ for $\vq=0.03$. A model that treats a particle with speed $0.5c$ by Newtonian dynamics cannot be taken as an approximation to the plasma at that time. The corresponding electron temperatures in the collisional study ($T_e/T_e(0)$ to $231$, $v_{te}>0.3c$ beyond a factor of $90$) leave the same model for the same reason. The nonrelativistic bounds of this chapter therefore remain valid for the model they are derived in, and for the early part of the evolution where the quiver speed is small (at $\vq=0.03$ and $\wpe t\lesssim200$ it is $\lesssim0.1c$), but they do not bound relativistic plasmas. Chapter~\ref{ch:rel} shows by how much and why the answer changes.

The same chapter shows one more point that bears on the interpretation of the windows. The positivity losses and the step-budget failures that end the certified windows of this chapter happen in the nonrelativistic solver at times, $800$--$1000$ for $\vq=0.03$ to $0.07$, at which the nonrelativistic law is already wrong in the relativistic models. The windows are therefore not limiting the nonrelativistic physics before it would have been replaced; they are limiting a numerical solution of a model that is no longer the right model. This is the sense in which the whole programme of Sections~\ref{sec:c9-windows} and~\ref{sec:c9-r1} was well posed and, for the HHS question, aimed at a model that turned out not to apply.
